\section{Don't Care Conditions}

\begin{definition}
A don't care condition occurs when the output of a logic circuit for a
specific input combination is irrelevant. These inputs are represented by \textit{d}
in Karnaugh-maps. Other representations that can be used are X, $\phi$
or $-$, these are commonly used in truth tables.
\end{definition}

\subsubsection{Why They Arise}

Don't care conditions typically arise in two situations:

\begin{description}
    \item[Invalid input combinations.] In a BCD (Binary Coded Decimal)
    system, only the input combinations \(0000\) to \(1001\) (representing
    digits 0--9) are valid. The remaining six combinations (\(1010\) to
    \(1111\)) can never occur, so their outputs are assigned as don't cares.

    \item[Unused outputs.] In some systems a particular input state is
    reachable, but the resulting output is never used by any other part of
    the circuit, making it irrelevant to the designer.
\end{description}

\subsubsection{How They Are Used}

During minimisation, don't care entries may be treated as either 0 or 1 —
whichever choice produces the simplest expression. This gives the designer
freedom to form \textbf{larger groupings (implicants)}, which directly
reduces the number of logic gates in the final circuit. A key rule applies:
a grouping may \emph{include} don't cares to enlarge it, but it cannot
consist \emph{entirely} of don't cares, as such a group covers no required
output.

\subsubsection{Truth Table Example}

Consider the function \(f(A,B,C) = \sum m(1,3,7) + \sum d(2,5)\), shown
in Table~\ref{tab:dontcare}. Minterms 2 and 5 are don't cares and can be
included in groupings during minimisation to yield a simpler final
expression.

\begin{table}[h!]
\centering
\caption{Truth table with don't care conditions}
\label{tab:dontcare}
\begin{tabular}{@{}ccccc@{}}
\toprule
\textbf{Minterm} & \textbf{A} & \textbf{B} & \textbf{C} &
\textbf{Output} \\
\midrule
0 & 0 & 0 & 0 & 0 \\
1 & 0 & 0 & 1 & 1 \\
2 & 0 & 1 & 0 & \(d\) \\
3 & 0 & 1 & 1 & 1 \\
4 & 1 & 0 & 0 & 0 \\
5 & 1 & 0 & 1 & \(d\) \\
6 & 1 & 1 & 0 & 0 \\
7 & 1 & 1 & 1 & 1 \\
\bottomrule
\end{tabular}
\end{table}

\section{Quine-McCluskey Method}

\begin{definition}
The Quine-McCluskey method is a tabular, algorithmic approach to Boolean
function minimisation. It is essentially a more systematic and
computer-friendly version of the Karnaugh map method. This method is a
deterministic algorithm that can be programmed and is effective for any
number of variables.
\end{definition}

\subsubsection{How It Works}

The Quine-McCluskey method proceeds in two well-defined stages.

\begin{description}
    \item[Step 1: Finding prime implicants.]
    This step identifies the largest possible groups of 1s and don't cares
    that can be combined. It works by first listing all minterms (and don't
    cares) in binary and sorting them into groups according to the number of
    1-bits in their binary representation. Terms from adjacent groups are
    then compared and merged if they differ in \textbf{exactly one bit
    position}, using the algebraic identity:
    \[
        AB + AB' = A
    \]
    The differing bit is replaced with a dash ($-$) to indicate elimination.
    This process repeats on the resulting terms until no further merging is
    possible. Any term that cannot be merged further is a
    \textbf{prime implicant}.

    \item[Step 2: Finding the essential prime implicants.]
    This step selects the minimal set of prime implicants needed to cover
    all the original minterms (the 1s of the function). A
    \textbf{prime implicant chart} is constructed with prime implicants as
    rows and original minterms as columns. A prime implicant is
    \textbf{essential} if it is the only one covering a particular minterm
    — it must be included in the final expression. After selecting all
    essential prime implicants, any uncovered minterms are resolved by
    choosing additional prime implicants, preferring those that cover the
    most remaining minterms. This is a classic set-covering problem.
\end{description}

\subsubsection{Worked Example}

Given: \(f(A,B,C,D) = \sum m(0,1,2,5,6,7,8,9,10,14)\)

All minterms are listed in binary and sorted by their number of 1-bits,
as shown in Table~\ref{tab:grouping}. Adjacent groups are then compared and
combined where possible, producing the first-round combinations in
Table~\ref{tab:combinations}.

\begin{table}[h!]
\centering
\caption{Minterms grouped by number of 1-bits}
\label{tab:grouping}
\begin{tabular}{@{}ccc@{}}
\toprule
\textbf{\# of 1s} & \textbf{Minterm} & \textbf{A\;B\;C\;D} \\
\midrule
0 & 0  & 0\;0\;0\;0 \\
\midrule
1 & 1  & 0\;0\;0\;1 \\
  & 2  & 0\;0\;1\;0 \\
  & 8  & 1\;0\;0\;0 \\
\midrule
2 & 5  & 0\;1\;0\;1 \\
  & 6  & 0\;1\;1\;0 \\
  & 9  & 1\;0\;0\;1 \\
  & 10 & 1\;0\;1\;0 \\
\midrule
3 & 7  & 0\;1\;1\;1 \\
  & 14 & 1\;1\;1\;0 \\
\bottomrule
\end{tabular}
\end{table}

\begin{table}[h!]
\centering
\caption{First-round combinations}
\label{tab:combinations}
\begin{tabular}{@{}lll@{}}
\toprule
\textbf{Minterms} & \textbf{Result} & \textbf{A\;B\;C\;D} \\
\midrule
0,\;1   & (0,1)  & 0\;0\;0\;-- \\
0,\;2   & (0,2)  & 0\;0\;--\;0 \\
0,\;8   & (0,8)  & --\;0\;0\;0 \\
1,\;9   & (1,9)  & --\;0\;0\;1 \\
2,\;6   & (2,6)  & 0\;--\;1\;0 \\
2,\;10  & (2,10) & --\;0\;1\;0 \\
8,\;9   & (8,9)  & 1\;0\;0\;-- \\
8,\;10  & (8,10) & 1\;0\;--\;0 \\
6,\;7   & (6,7)  & 0\;1\;1\;-- \\
6,\;14  & (6,14) & --\;1\;1\;0 \\
\bottomrule
\end{tabular}
\end{table}

Combining continues on these results until no further merging is possible.
Terms that remain uncombined are the prime implicants. These are entered
into the prime implicant chart, essential prime implicants are identified,
and the minimised Sum of Products expression is written.
